\section{Weiterentwicklung der Anwendung Mailbroadcast}
\label{sec:praxis}
Ziel dieses Kapitels ist es, 


\subsection{Mailbroadcast: Überblick und Einsatzbereiche}
wie in Abschnitt \ref{subsubsec:mbc_allianz} aufgegriffen


\cite{mbcSchulung}

\subsection{Ist-Zustandsanalyse der Anwendung Mailbroadcast}
Die Analyse des Weiterentwicklungsbedarfs umfasst primär die tiefgründige Untersuchung des Ist-Zustands der ... Anwendung Mailbroadcast, die ...

Schulung erwähnen

In 2.2 aus der Perspektive der Benutzer kurz vorgestellt hier Ist-Zustand umfassend analysiert

\subsubsection{Architektur}
\subsubsection{Bereitgestellte Personaldaten}
\subsubsection{Benutzeroberfläche in Microsoft Access}
\subsubsection{Bewertung der vorhandenen Dokumentation}
Weiterentwicklung umfasst auch eine gute Doku
\subsubsection{Zusammenfassung und identifizierte Schwachstellen}
Nur Daten behalten, die wir brauchen --> Daten von Ausgeschiedenen löschen; Index setzen zur Beschleunigung
Verzögerte Mailzustellung

Auf die Erfüllung der Anforderungen eingehen

bsp mit student in einem Referat aber eigentlich in der Berufsausbildung --> Problem bereits in Strukturverteilern in Outlook vorhanden
aber in MBC nicht explizit adressiert


\subsection{Projekt \enquote{Mailbroadcast Stufe 3}}

\subsubsection{Inhaltliche Erweiterung und fachliche Bereinigung der aktuellen Datenlieferung}
\subsubsection{Direktlieferung der Personaldaten von APAS-SAP} % ohne Host-Schnittstelle}
\subsubsection{Funktionale Erweiterungen}
\subsubsection{Umstellung der Datenlieferung auf SuccessFactors}